% GENERATED FILE. Edit canonical lessons, then run npm run generate:books.
% canonical-source-hash: fnv1a64:8dd59987cd3dbb0e
% canonical-lessons: RU-C375-garantirovat, RU-C375-obyazatelno, RU-C375-neobkhodimyy, RU-C375-neobkhodimost, RU-C375-dostatochno

\chapter{Without fail}
\label{ch:ru-without-fail}
\hlchaptermodality{375}

\begin{chapteropening}
\textbf{By the end of this chapter:} \emph{I can say that I will do something without fail, that something is necessary or is enough, and that I can guarantee it.}

Complete the last lesson of chapter 375: I can say that I will do something without fail, that something is necessary or is enough, and that I can guarantee it.
\end{chapteropening}

\section[\texorpdfstring{\ru{гарантировать} (garantírovat)}{гарантировать (garantírovat)}]{\ru{гарантировать} (garantírovat) — to guarantee}

\label{lesson:RU-C375-garantirovat}



\begin{quote}
Before the new one: say the Russian for to succeed, then the Russian for to work out.
\end{quote}

\subsection*{You'll want to know: \ru{гарантировать}}
\textbf{\ru{гарантировать}} (\emph{garantírovat}) — \textquotedblleft{}to guarantee\textquotedblright{}.

\begin{quote}
\ru{Я} \ru{не} \ru{могу} \ru{гарантировать}, \ru{что} \ru{всё} \ru{получится}.
\end{quote}

\begin{quote}
\emph{Ya ne mogú garantírovat, chto vsyo polúchitsya.}
\end{quote}

\begin{quote}
I cannot guarantee that everything will work out.
\end{quote}

The same form serves as both imperfective and perfective; present \emph{\ru{гарантирую}, \ru{гарантируешь}}.

\begin{cousinweb}
\textbf{\ru{гарантировать}} is usually traced, through German, to French \emph{garantir} \textquotedblleft{}to guarantee\textquotedblright{} (EtymDB), built on \emph{garant} \textquotedblleft{}a guarantor\textquotedblright{}; EtymDB sets \emph{garant} beside Old High German \emph{werento} \textquotedblleft{}a guarantor\textquotedblright{} (EtymDB).
\end{cousinweb}

\begin{grammarlens}[title={What you've built}]
One more word for telling, explaining and describing at B1.
\end{grammarlens}

\subsection*{Your turn}
\begin{itemize}
\raggedright
  \item \emph{Say it:} \emph{garantírovat}
  \item \emph{Say it:} \emph{garantírovat}, once more
  \item \emph{Say it:} a sentence of your own with \emph{garantírovat}
  \item \emph{Recall:} say the Russian for to succeed, then the Russian for to work out, then say \emph{garantírovat} again
\end{itemize}

\begin{tcolorbox}[breakable,colback=teal!4,colframe=teal!35!black,title={Before you move on}]
What does \ru{гарантировать} mean? (\textquotedblleft{}To guarantee\textquotedblright{}.) Say it once more.
\end{tcolorbox}
\section[\texorpdfstring{\ru{обязательно} (obyazátelno)}{обязательно (obyazátelno)}]{\ru{обязательно} (obyazátelno) — definitely, without fail}

\label{lesson:RU-C375-obyazatelno}



\begin{quote}
Before the new one: say the Russian for to work out, then the Russian for to guarantee.
\end{quote}

\subsection*{You'll want to know: \ru{обязательно}}
\textbf{\ru{обязательно}} (\emph{obyazátelno}) — \textquotedblleft{}definitely, without fail\textquotedblright{}.

\begin{quote}
\ru{Я} \ru{обязательно} \ru{позвоню} \ru{тебе} \ru{завтра}.
\end{quote}

\begin{quote}
\emph{Ya obyazátelno pozvonyú tebé závtra.}
\end{quote}

\begin{quote}
I will definitely call you tomorrow.
\end{quote}

It also means necessarily: \emph{\ru{не} \ru{обязательно}} — not necessarily.

\begin{cousinweb}
\textbf{\ru{обязательно}} is built on \emph{\ru{обязать}} \textquotedblleft{}to oblige\textquotedblright{}, from \emph{\ru{о}-} and \emph{\ru{вязать}} \textquotedblleft{}to tie, to bind\textquotedblright{}, Proto-Slavic \emph{*vęzati}: bound to happen.
\end{cousinweb}

\begin{grammarlens}[title={What you've built}]
One more word for telling, explaining and describing at B1.
\end{grammarlens}

\subsection*{Your turn}
\begin{itemize}
\raggedright
  \item \emph{Say it:} \emph{obyazátelno}
  \item \emph{Say it:} \emph{obyazátelno}, once more
  \item \emph{Say it:} a sentence of your own with \emph{obyazátelno}
  \item \emph{Recall:} say the Russian for to work out, then the Russian for to guarantee, then say \emph{obyazátelno} again
\end{itemize}

\begin{tcolorbox}[breakable,colback=teal!4,colframe=teal!35!black,title={Before you move on}]
What does \ru{обязательно} mean? (\textquotedblleft{}Definitely, without fail\textquotedblright{}.) Say it once more.
\end{tcolorbox}
\section[\texorpdfstring{\ru{необходимый} (neobkhodímyy)}{необходимый (neobkhodímyy)}]{\ru{необходимый} (neobkhodímyy) — necessary, essential}

\label{lesson:RU-C375-neobkhodimyy}



\begin{quote}
Before the new one: say the Russian for to guarantee, then the Russian for definitely.
\end{quote}

\subsection*{You'll want to know: \ru{необходимый}}
\textbf{\ru{необходимый}} (\emph{neobkhodímyy}) — \textquotedblleft{}necessary, essential\textquotedblright{}.

\begin{quote}
\ru{У} \ru{нас} \ru{есть} \ru{всё} \ru{необходимое}.
\end{quote}

\begin{quote}
\emph{U nas yest vsyo neobkhodímoye.}
\end{quote}

\begin{quote}
We have everything we need.
\end{quote}

\emph{\ru{Необходимо}} + infinitive — it is necessary to. More formal than \emph{\ru{нужный}}.

\begin{cousinweb}
\textbf{\ru{необходимый}} is built from \emph{\ru{не}-} \textquotedblleft{}not\textquotedblright{} and \emph{\ru{обходить}} \textquotedblleft{}to go round\textquotedblright{}, from \emph{\ru{ходить}} \textquotedblleft{}to go\textquotedblright{}, Proto-Slavic \emph{*xoditi}: what cannot be got round.
\end{cousinweb}

\begin{grammarlens}[title={What you've built}]
One more word for telling, explaining and describing at B1.
\end{grammarlens}

\subsection*{Your turn}
\begin{itemize}
\raggedright
  \item \emph{Say it:} \emph{neobkhodímyy}
  \item \emph{Say it:} \emph{neobkhodímyy}, once more
  \item \emph{Say it:} a sentence of your own with \emph{neobkhodímyy}
  \item \emph{Recall:} say the Russian for to guarantee, then the Russian for definitely, then say \emph{neobkhodímyy} again
\end{itemize}

\begin{tcolorbox}[breakable,colback=teal!4,colframe=teal!35!black,title={Before you move on}]
What does \ru{необходимый} mean? (\textquotedblleft{}Necessary, essential\textquotedblright{}.) Say it once more.
\end{tcolorbox}
\section[\texorpdfstring{\ru{необходимость} (neobkhodímost)}{необходимость (neobkhodímost)}]{\ru{необходимость} (neobkhodímost) — necessity, need}

\label{lesson:RU-C375-neobkhodimost}



\begin{quote}
Before the new one: say the Russian for definitely, then the Russian for necessary.
\end{quote}

\subsection*{You'll want to know: \ru{необходимость}}
\textbf{\ru{необходимость}} (\emph{neobkhodímost}) — \textquotedblleft{}necessity, need\textquotedblright{}.

\begin{quote}
\ru{В} \ru{этом} \ru{нет} \ru{необходимости}.
\end{quote}

\begin{quote}
\emph{V étom net neobkhodímosti.}
\end{quote}

\begin{quote}
There is no need for that.
\end{quote}

Feminine. \emph{\ru{В} \ru{случае} \ru{необходимости}} — if need be.

\begin{cousinweb}
\textbf{\ru{необходимость}} is built on \emph{\ru{необходимый}} \textquotedblleft{}necessary\textquotedblright{} with the abstract suffix \emph{-\ru{ость}}; the root is \emph{\ru{ходить}} \textquotedblleft{}to go\textquotedblright{}.
\end{cousinweb}

\begin{grammarlens}[title={What you've built}]
One more word for telling, explaining and describing at B1.
\end{grammarlens}

\subsection*{Your turn}
\begin{itemize}
\raggedright
  \item \emph{Say it:} \emph{neobkhodímost}
  \item \emph{Say it:} \emph{neobkhodímost}, once more
  \item \emph{Say it:} a sentence of your own with \emph{neobkhodímost}
  \item \emph{Recall:} say the Russian for definitely, then the Russian for necessary, then say \emph{neobkhodímost} again
\end{itemize}

\begin{tcolorbox}[breakable,colback=teal!4,colframe=teal!35!black,title={Before you move on}]
What does \ru{необходимость} mean? (\textquotedblleft{}Necessity, need\textquotedblright{}.) Say it once more.
\end{tcolorbox}
\section[\texorpdfstring{\ru{достаточно} (dostátochno)}{достаточно (dostátochno)}]{\ru{достаточно} (dostátochno) — enough; fairly}

\label{lesson:RU-C375-dostatochno}



\begin{quote}
Before the new one: say the Russian for necessary, then the Russian for necessity.
\end{quote}

\subsection*{You'll want to know: \ru{достаточно}}
\textbf{\ru{достаточно}} (\emph{dostátochno}) — \textquotedblleft{}enough; fairly\textquotedblright{}.

\begin{quote}
\ru{У} \ru{нас} \ru{достаточно} \ru{времени}.
\end{quote}

\begin{quote}
\emph{U nas dostátochno vrémeni.}
\end{quote}

\begin{quote}
We have enough time.
\end{quote}

With the genitive, enough: \emph{\ru{достаточно} \ru{денег}}. Before an adjective, fairly: \emph{\ru{достаточно} \ru{трудно}}.

\begin{cousinweb}
\textbf{\ru{достаточно}} is built on \emph{\ru{достаток}} \textquotedblleft{}sufficiency\textquotedblright{}, from \emph{\ru{достать}} \textquotedblleft{}to reach, to be enough\textquotedblright{}, on \emph{\ru{стать}}, the partner of \emph{\ru{стоять}}, Proto-Slavic \emph{*stojati} (EtymDB).
\end{cousinweb}

\begin{grammarlens}[title={What you've built}]
One more word for telling, explaining and describing at B1.
\end{grammarlens}

\subsection*{Your turn}
\begin{itemize}
\raggedright
  \item \emph{Say it:} \emph{dostátochno}
  \item \emph{Say it:} \emph{dostátochno}, once more
  \item \emph{Say it:} a sentence of your own with \emph{dostátochno}
  \item \emph{Recall:} say the Russian for necessary, then the Russian for necessity, then say \emph{dostátochno} again
\end{itemize}

\begin{tcolorbox}[breakable,colback=teal!4,colframe=teal!35!black,title={Before you move on}]
What does \ru{достаточно} mean? (\textquotedblleft{}Enough; fairly\textquotedblright{}.) Say it once more.
\end{tcolorbox}
